\documentclass[reqno]{amsart}
\usepackage[T1]{fontenc}
\usepackage{amsmath,amssymb,amsthm}
\usepackage{xr-hyper}
\usepackage[hidelinks]{hyperref}
\externaldocument[ec-]{20261124T223000--elliptic-curves-for-the-proof__workbook}
\externaldocument[kl-]{20261125T009000--kronecker-limit-and-chowla-selberg__workbook}
\externaldocument[cp-]{20261125T006500--covolumes-and-periods__workbook}
\externaldocument[hp-]{20261125T020300--heights-periods-and-the-auxiliary-prime__workbook}
\externaldocument[sg-]{20261126T020320--slopes-of-lattices-of-forms__workbook}
\externaldocument[lc-]{20261126T020340--lower-slope-bounds-on-covers__workbook}

\newcounter{chap}
% Chapter numbers, generated from 20261125T012000--grh-workbook__guide.org
% by scripts/make-chapters.py.  Do not edit: change the order of the
% index and rerun the script (./preview-site.sh and CI do this).
%
% Usage in a sheet, after \newcounter{chap}:
%   % Chapter numbers, generated from 20261125T012000--grh-workbook__guide.org
% by scripts/make-chapters.py.  Do not edit: change the order of the
% index and rerun the script (./preview-site.sh and CI do this).
%
% Usage in a sheet, after \newcounter{chap}:
%   % Chapter numbers, generated from 20261125T012000--grh-workbook__guide.org
% by scripts/make-chapters.py.  Do not edit: change the order of the
% index and rerun the script (./preview-site.sh and CI do this).
%
% Usage in a sheet, after \newcounter{chap}:
%   \input{chapters}\chapnumber{<sheet ID>}
\makeatletter
\providecommand\chapters@set[2]{\@namedef{chapters@#1}{#2}}
\chapters@set{20261125T001000}{1}% zeros-to-ramp-sums
\chapters@set{20261125T001100}{2}% inert-primes-and-estimate-b
\chapters@set{20261125T005000}{3}% galois-fourier-and-averaging
\chapters@set{20261125T005500}{4}% columns-from-a-sphere
\chapters@set{20261125T000000}{5}% slope-method
\chapters@set{20261126T030100}{6}% a-twist-for-each-prime
\chapters@set{20261125T003000}{7}% the-chain-and-the-square-root
\chapters@set{20261125T010000}{8}% rank-lemma-and-zero-estimate
\chapters@set{20261125T010500}{9}% poincare-series-and-near-orthogonality
\chapters@set{20261124T223000}{10}% elliptic-curves-for-the-proof
\chapters@set{20261125T009500}{11}% fock-model-and-theta-intertwiners
\chapters@set{20261124T223100}{12}% complex-multiplication-and-inert-reduction
\chapters@set{20261125T008000}{13}% one-cm-determinant
\chapters@set{20261124T214500}{14}% supersingular-theta-divisibility
\chapters@set{20261124T214600}{15}% the-isotropic-plane
\chapters@set{20261125T013700}{16}% pairings-on-a-formal-group-of-height-two
\chapters@set{20261125T013500}{17}% hecke-nondegeneracy-on-the-tree
\chapters@set{20261126T030200}{18}% the-bound-is-attained
\chapters@set{20261125T009000}{19}% kronecker-limit-and-chowla-selberg
\chapters@set{20261125T002000}{20}% definite-theta-lifts
\chapters@set{20261125T015500}{21}% siegel-weil-in-the-smallest-case
\chapters@set{20261125T002100}{22}% definite-theta-lifts-and-norms
\chapters@set{20261125T016500}{23}% hecke-operators-and-theta-lifts
\chapters@set{20261125T016700}{24}% divisibility-along-the-chain
\chapters@set{20261125T004000}{25}% denominators-and-sha
\chapters@set{20261125T004100}{26}% tate-twists-and-sha
\chapters@set{20261126T020430}{27}% the-minus-part-by-gauss-sums
\chapters@set{20261125T014100}{28}% congruence-modules
\chapters@set{20261125T014500}{29}% eisenstein-congruences-and-selmer-classes
\chapters@set{20261125T014700}{30}% pseudocharacters-and-the-length-bound
\chapters@set{20261125T011000}{31}% galois-representations-black-box
\chapters@set{20261126T030300}{32}% the-data
\chapters@set{20261126T020210}{33}% bessel-poincare-series-and-their-gram-operators
\chapters@set{20261126T030400}{34}% the-seesaw-identity
\chapters@set{20261126T030500}{35}% divisibility-of-the-composite-lift
\chapters@set{20261126T030600}{36}% theorem-a-short-forms
\chapters@set{20261126T030700}{37}% the-determinant-of-a-chain
\chapters@set{20261126T030800}{38}% the-cost-of-one-lift
\chapters@set{20261126T030900}{39}% denominators-of-projections-to-lifts
\chapters@set{20261126T031000}{40}% twists-with-few-congruences
\chapters@set{20261126T031100}{41}% theorem-c-and-the-main-inequality
\providecommand\chapnumber[1]{%
  \@ifundefined{chapters@#1}%
    {\PackageError{chapters}{Sheet #1 is not listed in the index}%
      {List it in 20261125T012000--grh-workbook__guide.org\MessageBreak
       and run scripts/make-chapters.py.}}%
    {\edef\thechap{\@nameuse{chapters@#1}}}}
\makeatother
\chapnumber{<sheet ID>}
\makeatletter
\providecommand\chapters@set[2]{\@namedef{chapters@#1}{#2}}
\chapters@set{20261125T001000}{1}% zeros-to-ramp-sums
\chapters@set{20261125T001100}{2}% inert-primes-and-estimate-b
\chapters@set{20261125T005000}{3}% galois-fourier-and-averaging
\chapters@set{20261125T005500}{4}% columns-from-a-sphere
\chapters@set{20261125T000000}{5}% slope-method
\chapters@set{20261126T030100}{6}% a-twist-for-each-prime
\chapters@set{20261125T003000}{7}% the-chain-and-the-square-root
\chapters@set{20261125T010000}{8}% rank-lemma-and-zero-estimate
\chapters@set{20261125T010500}{9}% poincare-series-and-near-orthogonality
\chapters@set{20261124T223000}{10}% elliptic-curves-for-the-proof
\chapters@set{20261125T009500}{11}% fock-model-and-theta-intertwiners
\chapters@set{20261124T223100}{12}% complex-multiplication-and-inert-reduction
\chapters@set{20261125T008000}{13}% one-cm-determinant
\chapters@set{20261124T214500}{14}% supersingular-theta-divisibility
\chapters@set{20261124T214600}{15}% the-isotropic-plane
\chapters@set{20261125T013700}{16}% pairings-on-a-formal-group-of-height-two
\chapters@set{20261125T013500}{17}% hecke-nondegeneracy-on-the-tree
\chapters@set{20261126T030200}{18}% the-bound-is-attained
\chapters@set{20261125T009000}{19}% kronecker-limit-and-chowla-selberg
\chapters@set{20261125T002000}{20}% definite-theta-lifts
\chapters@set{20261125T015500}{21}% siegel-weil-in-the-smallest-case
\chapters@set{20261125T002100}{22}% definite-theta-lifts-and-norms
\chapters@set{20261125T016500}{23}% hecke-operators-and-theta-lifts
\chapters@set{20261125T016700}{24}% divisibility-along-the-chain
\chapters@set{20261125T004000}{25}% denominators-and-sha
\chapters@set{20261125T004100}{26}% tate-twists-and-sha
\chapters@set{20261126T020430}{27}% the-minus-part-by-gauss-sums
\chapters@set{20261125T014100}{28}% congruence-modules
\chapters@set{20261125T014500}{29}% eisenstein-congruences-and-selmer-classes
\chapters@set{20261125T014700}{30}% pseudocharacters-and-the-length-bound
\chapters@set{20261125T011000}{31}% galois-representations-black-box
\chapters@set{20261126T030300}{32}% the-data
\chapters@set{20261126T020210}{33}% bessel-poincare-series-and-their-gram-operators
\chapters@set{20261126T030400}{34}% the-seesaw-identity
\chapters@set{20261126T030500}{35}% divisibility-of-the-composite-lift
\chapters@set{20261126T030600}{36}% theorem-a-short-forms
\chapters@set{20261126T030700}{37}% the-determinant-of-a-chain
\chapters@set{20261126T030800}{38}% the-cost-of-one-lift
\chapters@set{20261126T030900}{39}% denominators-of-projections-to-lifts
\chapters@set{20261126T031000}{40}% twists-with-few-congruences
\chapters@set{20261126T031100}{41}% theorem-c-and-the-main-inequality
\providecommand\chapnumber[1]{%
  \@ifundefined{chapters@#1}%
    {\PackageError{chapters}{Sheet #1 is not listed in the index}%
      {List it in 20261125T012000--grh-workbook__guide.org\MessageBreak
       and run scripts/make-chapters.py.}}%
    {\edef\thechap{\@nameuse{chapters@#1}}}}
\makeatother
\chapnumber{<sheet ID>}
\makeatletter
\providecommand\chapters@set[2]{\@namedef{chapters@#1}{#2}}
\chapters@set{20261125T001000}{1}% zeros-to-ramp-sums
\chapters@set{20261125T001100}{2}% inert-primes-and-estimate-b
\chapters@set{20261125T005000}{3}% galois-fourier-and-averaging
\chapters@set{20261125T005500}{4}% columns-from-a-sphere
\chapters@set{20261125T000000}{5}% slope-method
\chapters@set{20261126T030100}{6}% a-twist-for-each-prime
\chapters@set{20261125T003000}{7}% the-chain-and-the-square-root
\chapters@set{20261125T010000}{8}% rank-lemma-and-zero-estimate
\chapters@set{20261125T010500}{9}% poincare-series-and-near-orthogonality
\chapters@set{20261124T223000}{10}% elliptic-curves-for-the-proof
\chapters@set{20261125T009500}{11}% fock-model-and-theta-intertwiners
\chapters@set{20261124T223100}{12}% complex-multiplication-and-inert-reduction
\chapters@set{20261125T008000}{13}% one-cm-determinant
\chapters@set{20261124T214500}{14}% supersingular-theta-divisibility
\chapters@set{20261124T214600}{15}% the-isotropic-plane
\chapters@set{20261125T013700}{16}% pairings-on-a-formal-group-of-height-two
\chapters@set{20261125T013500}{17}% hecke-nondegeneracy-on-the-tree
\chapters@set{20261126T030200}{18}% the-bound-is-attained
\chapters@set{20261125T009000}{19}% kronecker-limit-and-chowla-selberg
\chapters@set{20261125T002000}{20}% definite-theta-lifts
\chapters@set{20261125T015500}{21}% siegel-weil-in-the-smallest-case
\chapters@set{20261125T002100}{22}% definite-theta-lifts-and-norms
\chapters@set{20261125T016500}{23}% hecke-operators-and-theta-lifts
\chapters@set{20261125T016700}{24}% divisibility-along-the-chain
\chapters@set{20261125T004000}{25}% denominators-and-sha
\chapters@set{20261125T004100}{26}% tate-twists-and-sha
\chapters@set{20261126T020430}{27}% the-minus-part-by-gauss-sums
\chapters@set{20261125T014100}{28}% congruence-modules
\chapters@set{20261125T014500}{29}% eisenstein-congruences-and-selmer-classes
\chapters@set{20261125T014700}{30}% pseudocharacters-and-the-length-bound
\chapters@set{20261125T011000}{31}% galois-representations-black-box
\chapters@set{20261126T030300}{32}% the-data
\chapters@set{20261126T020210}{33}% bessel-poincare-series-and-their-gram-operators
\chapters@set{20261126T030400}{34}% the-seesaw-identity
\chapters@set{20261126T030500}{35}% divisibility-of-the-composite-lift
\chapters@set{20261126T030600}{36}% theorem-a-short-forms
\chapters@set{20261126T030700}{37}% the-determinant-of-a-chain
\chapters@set{20261126T030800}{38}% the-cost-of-one-lift
\chapters@set{20261126T030900}{39}% denominators-of-projections-to-lifts
\chapters@set{20261126T031000}{40}% twists-with-few-congruences
\chapters@set{20261126T031100}{41}% theorem-c-and-the-main-inequality
\providecommand\chapnumber[1]{%
  \@ifundefined{chapters@#1}%
    {\PackageError{chapters}{Sheet #1 is not listed in the index}%
      {List it in 20261125T012000--grh-workbook__guide.org\MessageBreak
       and run scripts/make-chapters.py.}}%
    {\edef\thechap{\@nameuse{chapters@#1}}}}
\makeatother

\ifcsname chapters@20261126T020330\endcsname
  \chapnumber{20261126T020330}
\else
  \renewcommand{\thechap}{?}
\fi
\renewcommand{\thesection}{\thechap.\arabic{section}}
\theoremstyle{definition}
\newtheorem{problem}{Problem}[section]
\theoremstyle{plain}
\newtheorem{theorem}{Theorem}[section]
\theoremstyle{remark}
\newtheorem*{leanremark}{In the Lean formalization}
\newcommand{\Z}{\mathbb Z}
\newcommand{\Q}{\mathbb Q}
\newcommand{\C}{\mathbb C}
\newcommand{\Oc}{\mathcal O}
\newcommand{\Ac}{\mathcal A}
\newcommand{\aaa}{\mathfrak a}
\newcommand{\ii}{\mathrm i}
\newcommand{\rd}{\operatorname{rd}}
\newcommand{\mus}{\mu^\sharp}
\newcommand{\hF}{h_{\mathrm F}}
\newcommand{\tl}[1]{\texttt{#1}}
\setlength{\emergencystretch}{3em}

\title{The CM period and the flat line}
\author{}
\date{November 26, 2026}

\begin{document}
\maketitle
\refstepcounter{chap}\label{chap}

\paragraph{Prerequisites.}
From Chapter~\ref{kl-chap}, \emph{The Kronecker limit formula and Chowla--Selberg}: the reduced forms \(Q\), the points \(\tau_Q\) with \(y_Q = \operatorname{Im}\tau_Q\), Problems~\ref{kl-pr:lerch-gen} and~\ref{kl-pr:LL-eta-gen}, all in Section~\ref{kl-sec:general}.  From Chapter~\ref{ec-chap}: Problems~\ref{ec-pr:varpi} and~\ref{ec-pr:Omega}, and \(\Delta(\Z + \Z\tau) = (2\pi)^{12}\eta(\tau)^{24}\).  From Chapter~\ref{cp-chap}, \emph{Covolumes and periods}: the Faltings height \(\hF\) in Deligne's normalization, Problems~\ref{cp-pr:homog} and~\ref{cp-pr:hF-i}, and Problem~\ref{cp-pr:dR-shift}.  From Chapter~\ref{hp-chap}: Theorem~\ref{hp-in:period} and its notation, and Problem~\ref{hp-pr:gaussian}.  From Chapter~\ref{sg-chap}: \(\mus\) and the main term \(T_j\) of Section~\ref{sg-sec:reduction}.  The theory of complex multiplication for elliptic curves is used through three classical facts, recalled in Section~\ref{sec:constant}.

The covolume estimate of Theorem~\ref{hp-in:A4} has the main term \(T_j = \sum_\rho(p_j - q_j)w_j(\rho)\log\Omega\).  Which period is \(\Omega\), and how does it reach a lattice of automorphic forms on a unitary Shimura variety?  In the Lean formalization \cite{Lean} the answer involves only elliptic curves: \(\Omega\) is the Chowla--Selberg period of an elliptic curve with complex multiplication by \(\Oc_K\), \(K = \Q(\sqrt{-d})\), and it enters through the degree of one metrized line, the \emph{flat line}, at points of the Shimura variety built from two such curves.

\section{The constant}\label{sec:constant}

Let \(d > 4\) be prime with \(d \equiv 3 \pmod 4\), \(K = \Q(\sqrt{-d})\) of class number \(h\), \(\chi\) its character, and \(c_\chi = L'(1,\chi)/L(1,\chi)\).  For a lattice \(L \subset \C\) put \(\Ac(L) := A(L)\lvert\Delta(L)\rvert^{1/6}\), which is invariant under homothety (Problem~\ref{cp-pr:homog}).  Let
\[
  \log\Omega_K := \gamma + \log\frac{2\pi}{\sqrt d} - c_\chi .
\]
For a CM field \(E \supseteq K\) as in Theorem~\ref{hp-in:period}, the constant there is \(\log\Omega = \log\Omega_K - \log(\rd_E/\sqrt d)\); for \(E = K\) the two agree.

\begin{problem}\label{pr:area-eta}
\(\Ac(\Z + \Z\tau) = 4\pi^2\operatorname{Im}(\tau)\,\lvert\eta(\tau)\rvert^4\).
\end{problem}

\begin{problem}\label{pr:class-mean}
\(\displaystyle \log\Omega_K = \frac1h\sum_Q\log\frac{\Ac(\Z + \Z\tau_Q)}{\pi} = \frac1h\sum_Q\log\bigl(4\pi y_Q\lvert\eta(\tau_Q)\rvert^4\bigr)\).
\end{problem}

\begin{problem}\label{pr:values}
\(\log\Omega_K = 1.426502\ldots\), \(1.300363\ldots\), \(1.027786\ldots\) for \(d = 7, 11, 19\).  For \(d = 23\) the three classes give \(\log(4\pi y_Q\lvert\eta(\tau_Q)\rvert^4) = 0.894533\ldots\) for \(Q = [1,1,6]\) and \(1.456932\ldots\) for \(Q = [2,\pm1,3]\), with mean \(\log\Omega_K = 1.269466\ldots\)
\end{problem}

The formula of Problem~\ref{pr:class-mean} also holds for \(d = 4\), where it gives \(1.476335\ldots\) (Problem~\ref{hp-pr:gaussian}).

\begin{problem}\label{pr:rd-E}
If \(E \supseteq K\) is a number field, \(\rd_E := \lvert\operatorname{disc}E\rvert^{1/[E:\Q]}\) (for CM fields this is the root discriminant of Chapter~\ref{hp-chap}) and \(\mathfrak d_{E/K}\) the relative discriminant, then \(\rd_E = \sqrt d\,N(\mathfrak d_{E/K})^{1/[E:\Q]}\).  Hence \(\log\Omega \le \log\Omega_K\), with equality exactly when \(E/K\) is unramified at every finite prime.
\end{problem}

The formalization does not average over ideal classes but over Galois conjugates of one lattice.  We use three classical facts.  First, the reduced form \(Q = [a, b, c]\) corresponds to the ideal \(a\Z + \frac{-b + \sqrt{-d}}{2}\Z\) of \(\Oc_K\), which is homothetic to \(\Z + \Z\tau_Q\), and this gives a bijection from reduced forms to ideal classes \cite[Thm.~7.7]{Co}.  Second, two lattices with the same \(j\) are homothetic \cite[Ch.~VI]{Si1}.  Third, the values \(j(\aaa)\), for \(\aaa\) running over representatives of the \(h\) ideal classes of \(K\), are distinct algebraic integers, and they are the roots of a monic irreducible polynomial over \(\Z\) of degree \(h\) \cite[Ch.~II]{Si2}.

\begin{problem}\label{pr:hilbert23}
For \(d = 23\), the three numbers \(j(\tau_Q)\) are \(-3493225.6999\ldots\) and \(737.8499\ldots \pm 1764.0189\ldots\ii\), and they are the roots of
\[
  X^3 + 3491750X^2 - 5151296875X + 12771880859375 .
\]
\end{problem}

\begin{problem}\label{pr:galois-mean}
Let \(L\) be a lattice homothetic to an ideal of \(\Oc_K\) with \(\Delta(L) = 1\), and let \(M \subset \C\) be a number field containing \(g_2(L)\) and \(g_3(L)\).  For each embedding \(\sigma \colon M \to \C\) let \(L^\sigma\) be the lattice with invariants \(\sigma g_2(L)\), \(\sigma g_3(L)\); these are the period lattices of the curve \(E_0 \colon y^2 = 4x^3 - g_2(L)x - g_3(L)\) over \(M\) with the differential \(dx/y\).  Then
\[
  \frac{1}{[M:\Q]}\sum_\sigma\log\frac{A(L^\sigma)}{\pi} = \log\Omega_K .
\]
\end{problem}

\begin{problem}\label{pr:conjugate}
For every lattice \(L\), the complex conjugate lattice \(\bar L\) has \(A(\bar L) = A(L)\) and \(\lvert\Delta(\bar L)\rvert = \lvert\Delta(L)\rvert\).  If \(L\) is homothetic to an ideal \(\aaa\), then \(\bar L\) is homothetic to \(\bar\aaa\), has \(\Delta(\bar L) = 1\) and invariants \(\overline{g_2(L)}\), \(\overline{g_3(L)}\).  So Problem~\ref{pr:galois-mean} applies to \(\bar L\) and the field \(\bar M\), and the conjugate curve \(\bar E_0\) has the same mean \(\log\Omega_K\).
\end{problem}

\emph{Status}: Problem~\ref{pr:galois-mean} is proved in the formalization from Chowla--Selberg in the form of Problem~\ref{pr:class-mean}, itself obtained there from Kronecker's limit formula.

\begin{problem}\label{pr:faltings}
If \(E_0\) is an elliptic curve over a number field with complex multiplication by \(\Oc_K\), then \(\hF(E_0) = -\frac12(\log\Omega_K + \log\pi)\).  For \(d = 7\) this is \(-1.285615\ldots\).  For \(K = \Q(\ii)\), where \(h = 1\) and the formula of Problem~\ref{pr:class-mean} also holds, the same argument gives \(-1.310532\ldots\), in agreement with Problem~\ref{cp-pr:hF-i}.
\end{problem}

\section{The main term}\label{sec:main}

The \(j\)th lift lives on a unitary group of signature \((p_j, q_j) = (k+j-1, m-j)\), \(1 \le j \le m\), so \(p_j + q_j = k + m - 1\) does not depend on \(j\).  In the formalization \(k \ge m + 3\).

\begin{problem}\label{pr:signature}
\(p_j - q_j = k + 2j - 1 - m\), and \(\sum_{j=1}^m(p_j - q_j) = mk\).  If \(k \ge m + 3\), then \(p_j - q_j \ge 4\) for every \(j\), and \(q_m = 0\).
\end{problem}

So only the last lift lives on a definite group, and every main term \(T_j\) has a positive coefficient in front of each weight.  The slope moves linearly in the weight when a lattice of sections is tensored with a line.

\begin{problem}\label{pr:tensor-line}
Let \(M\) be a nonzero lattice and \(\ell\) a lattice of rank one, in normed spaces over a number field, and give \(M \otimes \ell\) the tensor product norms and unit lattices.  Then \(\mus(M \otimes \ell) = \mus(M) + \mus(\ell)\).  Hence \(\mus(M \otimes \ell^{\otimes w}) = \mus(M) + w\,\mus(\ell)\) for \(w \ge 0\).
\end{problem}

\begin{problem}\label{pr:definite-shift}
For the lemniscatic lattice \(\varpi\Z[\ii]\), \(\Ac(\varpi\Z[\ii])/\pi = 2\varpi^2/\pi\).  So for \(K = \Q(\ii)\) the constant \(\Omega\) of Theorem~\ref{hp-in:period} is the number \(\Omega\) of Problem~\ref{ec-pr:Omega}, and the shift \(\lvert\lambda\rvert\log A(c\,\Oc)\) of Problem~\ref{cp-pr:dR-shift} equals \(\lvert\lambda\rvert(\log\Omega + \log\frac\pi2)\).
\end{problem}

Thus on a definite group (\(q = 0\)) the period enters the slope of the lattice of forms linearly in the weight, with coefficient \(\log\Omega\) up to a constant per unit of weight.  In the formalization the same happens in signature \((p, q)\), with coefficient \((p-q)\log\Omega\) per unit of weight, through the flat line of the next section.

\section{The flat line}\label{sec:flat}

Let \(X\) be the unitary Shimura variety of signature \((p, q)\) attached to \(E \supseteq K\), and let \(E_0\) and \(E_0'\) be elliptic curves with complex multiplication by \(\Oc_K\) through the two embeddings of \(K\), with \(\Delta = 1\).  In the formalization certain points of \(X\) are built from \(E_0\) and \(E_0'\), with \(p\) and \(q\) copies respectively, and carry a line of invariant differentials of a fixed type; a section of this line that is constant along such points is a \emph{flat generator}, and its metrized line the \emph{flat line}.  We do not reproduce the definitions of the formalization, which fix the type and the metric through a Hermitian form on the \(\Oc_E\)-lattice of the point.

\begin{theorem}[flat degree]\label{thm:flat}
The flat line at a point of type \((p, q)\), over a field of definition \(F\), has degree divided by \([F:\Q]\) (half the normalization of Theorem~\ref{hp-in:period}) equal to
\[
  \frac{p - q}{2}\log\Omega + q\,\bigl(\log d - \log\rd_E\bigr) .
\]
\end{theorem}
\noindent\emph{Status}: proved in the formalization.  Its degree is a sum of local terms.  At the archimedean places it is the logarithm of a Gram determinant of CM vectors, expressed through the norms of the invariant differentials of \(E_0\) and \(E_0'\), whose mean over conjugates is \(\log\Omega_K\) (Problems~\ref{pr:galois-mean} and~\ref{pr:conjugate}).  At the finite places it is a transition determinant of absolute value \(1\) together with the different of \(\Oc_E\).  We have read this proof but not the definitions of the flat generator and of its metric, so we state only the result.

\begin{problem}\label{pr:flat-K}
For \(E = K\), the flat degree is \(\frac{p-q}2\log\Omega_K + \frac q2\log d\).
\end{problem}

\begin{problem}\label{pr:flat-lifts}
Let \(F_j\) be the flat degree in signature \((p_j, q_j)\).  Then \(F_{j+1} - F_j = \log\Omega_K - \frac12\log d\) for \(1 \le j < m\), whatever \(E\) is.  For \(d = 7\) this difference is \(0.453546\ldots\)
\end{problem}

\begin{problem}\label{pr:period-part}
Write \(F^\circ_j := \frac{p_j - q_j}{2}\log\Omega\) for the period part of \(F_j\).  Then \(T_j = 2\sum_\rho w_j(\rho)F^\circ_j\), and \(\sum_jT_j = fmkw\log\Omega\) when \(w_j(\rho) = w\) for all \(j\) and \(\rho\), where \(f\) is the number of places \(\rho \in \Phi\).
\end{problem}

By Problem~\ref{pr:tensor-line}, tensoring with \(w\) copies of a line whose \(\mus\) is twice its normalized degree \(F^\circ_j\) moves \(\mus\) by \(2wF^\circ_j\); Problem~\ref{pr:period-part} says that \(T_j\) is exactly this shift for the weights \(w_j(\rho)\).  The formalization uses this as follows.  An integer \(\alpha_\rho\), the \emph{anchor}, is attached to each place, with
\[
  \Bigl\lvert T_j - \sum_\rho(p_j - q_j)\,\alpha_\rho\log\Omega\Bigr\rvert \le Cu ,
\]
in the notation of Section~\ref{sg-sec:reduction}.  Multiplication by a flat section of the corresponding tensor power of the flat line then compares two weights on a cover.

\begin{problem}\label{pr:anchor-error}
If \(\lvert\alpha_\rho - w_j(\rho)\rvert \le \eta\) for every \(\rho\), then \(\bigl\lvert T_j - \sum_\rho(p_j - q_j)\alpha_\rho\log\Omega\bigr\rvert \le f(k + m - 1)\,\eta\,\lvert\log\Omega\rvert\).
\end{problem}

\begin{theorem}[flat comparison]\label{thm:flatat}
On a good cover of the Shimura variety of the \(j\)th lift, \(j < m\), the lattices of integral cusp forms of the weights \(\lambda_-\) and \(\lambda_M\) of Chapter~\ref{lc-chap} satisfy
\[
  \mus(\lambda_-) \ge \mus(\lambda_M) - T_j - C_Fu ,
\]
with \(C_F\) independent of the point and the lift.
\end{theorem}
\noindent\emph{Status}: proved in the formalization from Theorem~\ref{thm:flat}, the anchor bound and the statements on the cover listed at the end of Chapter~\ref{lc-chap}.  The anchor bound is proved there from the definitions of the anchors, which we have not traced; we do not know how close \(\alpha_\rho\) is to \(w_j(\rho)\).

\paragraph{Where else the constant occurs.}  The same closed formula \(\gamma + \log(2\pi/\rd_E) - c_\chi\) is used in the formalization for the main term \(T_j\), for the flat degree, and in the bounds for the theta kernels of the chain, where Problem~\ref{pr:galois-mean} is applied.  So the cancellation of the period between the covolumes and the heights of the period maps is built into the definitions.  Only elliptic curves with complex multiplication by \(\Oc_K\) enter this step; no period relation for abelian varieties with complex multiplication by a larger field is used.

\begin{leanremark}
\raggedright In \cite{Lean} \(\log\Omega\) is \tl{NumDatum.\allowbreak{}logOm} in the main term \tl{l3m} and \tl{GlobalTheory.\allowbreak{}CM.\allowbreak{}logOm} in the flat degree; \tl{logOmVal\_\allowbreak{}eq} identifies the value.  Problem~\ref{pr:class-mean} is \tl{chowla\_\allowbreak{}selberg\_\allowbreak{}eta} and \tl{chowla\_\allowbreak{}selberg\_\allowbreak{}lattice}; Problem~\ref{pr:galois-mean} is \tl{mean\_\allowbreak{}log\_\allowbreak{}period} and \tl{MeanPeriod} (proved as \tl{MeanPeriodClosed.\allowbreak{}meanPeriod}, used in \tl{factors\_\allowbreak{}leS}).  The curves \(E_0\), \(E_0'\) come from \tl{exists\_\allowbreak{}curve0\_\allowbreak{}hasPeriod}.  Theorem~\ref{thm:flat} is \tl{GlobalTheory.\allowbreak{}CompanionTheory.\allowbreak{}Ar.\allowbreak{}Fl.\allowbreak{}flatDegree}, assembled by \tl{deg\_\allowbreak{}eq\_\allowbreak{}of\_\allowbreak{}local\_\allowbreak{}two\_\allowbreak{}curves}; it reaches the covolume bound through \tl{flatDegreeSt\_\allowbreak{}of\_\allowbreak{}resK\_\allowbreak{}gaga} and \tl{flatDeg\_\allowbreak{}alt\_\allowbreak{}alt}.  The anchor bound is \tl{AnchorAt}, proved by \tl{anchor\_\allowbreak{}holds\_\allowbreak{}alt}, and Theorem~\ref{thm:flatat} is \tl{FlatAt}, proved by \tl{flat\_\allowbreak{}alt\_\allowbreak{}alt}.
\end{leanremark}

\section*{Hints}

\subsection*{\S\ref{sec:constant}}
\begin{description}
\item[\ref{pr:area-eta}] \(A(\Z + \Z\tau) = \operatorname{Im}\tau\).
\item[\ref{pr:class-mean}] Problem~\ref{kl-pr:LL-eta-gen} and Problem~\ref{pr:area-eta}.
\item[\ref{pr:values}] Problem~\ref{pr:class-mean}, or Problem~\ref{kl-pr:lerch-gen} for \(c_\chi\).
\item[\ref{pr:rd-E}] The discriminant tower formula \(\lvert\operatorname{disc}E\rvert = \lvert\operatorname{disc}K\rvert^{[E:K]}N(\mathfrak d_{E/K})\).
\item[\ref{pr:hilbert23}] The \(\tau_Q\) of the three reduced forms of Problem~\ref{pr:values} are as in Section~\ref{kl-sec:general}; compute \(j = E_4^3/\Delta\) numerically, and round the symmetric functions, which are integers.
\item[\ref{pr:galois-mean}] \(\Ac(L^\sigma) = A(L^\sigma)\) depends only on \(j(L^\sigma) = \sigma j(L)\); count the embeddings of \(M\) above each embedding of \(\Q(j(L))\).
\item[\ref{pr:faltings}] The second formula for \(\hF\) in Chapter~\ref{cp-chap}, which does not depend on the model; \(\Ac\) is homothety invariant (Problem~\ref{cp-pr:homog}), so Problem~\ref{pr:galois-mean} applies after rescaling.
\end{description}

\subsection*{\S\ref{sec:main}}
\begin{description}
\item[\ref{pr:tensor-line}] A generator of \(\det(M \otimes \ell)\) is \(\omega \otimes e^{\otimes r}\) for generators \(\omega\) of \(\det M\) and \(e\) of \(\ell\).
\item[\ref{pr:definite-shift}] Problem~\ref{ec-pr:varpi} gives \(\Delta(\varpi\Z[\ii]) = 64\).
\end{description}

\subsection*{\S\ref{sec:flat}}
\begin{description}
\item[\ref{pr:flat-lifts}] Problem~\ref{pr:signature}.
\item[\ref{pr:period-part}] Problem~\ref{pr:signature}.
\end{description}

\begin{thebibliography}{Lean}
\bibitem[Lean]{Lean} The Lean community, \emph{riemann-hypothesis}, a Lean formalization of Alp\"oge's argument, \url{https://github.com/leanprover-community/riemann-hypothesis}.
\bibitem[Co]{Co} D.~A.~Cox, \emph{Primes of the Form \(x^2 + ny^2\)}, 2nd ed., Wiley, 2013.
\bibitem[Si1]{Si1} J.~H.~Silverman, \emph{The Arithmetic of Elliptic Curves}, 2nd ed., Grad.\ Texts in Math.\ 106, Springer, 2009.
\bibitem[Si2]{Si2} J.~H.~Silverman, \emph{Advanced Topics in the Arithmetic of Elliptic Curves}, Grad.\ Texts in Math.\ 151, Springer, 1994.
\end{thebibliography}
\end{document}
